\textbf{Phần 1. Trắc nghiệm lựa chọn}
\setcounterpageref{ex}{0}

\begin{ex} %[3D1V47]
Chữ số La Mã $V$ biểu thị số tự nhiên nào?
\choice
{\True $5$}
{$1$}
{$10$}
{$15$}
\loigiai{
Chữ số La Mã V biểu thị số 5.
\textbf{Đáp án của câu hỏi là A}
}
\end{ex}

\begin{ex} %[3D1V47]
Số 9 được viết bằng chữ số La Mã là
\choice
{$\text{VIIII}$}
{$\text{XI}$}
{\True $\text{IX}$}
{$\text{VI}$}
\loigiai{
Số 9 viết là IX (10 bớt 1).
\textbf{Đáp án của câu hỏi là C}
}
\end{ex}

\begin{ex} %[3D1V47]
Chữ số La Mã $\text{XIV}$ có giá trị trong hệ thập phân là
\choice
{$16$}
{\True $14$}
{$15$}
{$13$}
\loigiai{
$\text{XIV} = 10 + (5 - 1) = 14$.
\textbf{Đáp án của câu hỏi là B}
}
\end{ex}

\begin{ex} %[3D1V47]
Đồng hồ chỉ kim ngắn ở số $\text{VIII}$ và kim dài ở số $\text{XII}$. Khi đó đồng hồ chỉ mấy giờ?
\choice
{$7$ giờ đúng}
{$9$ giờ đúng}
{$12$ giờ đúng}
{\True $8$ giờ đúng}
\loigiai{
$\text{VIII}$ là số 8, $\text{XII}$ là số 12. Đồng hồ chỉ đúng 8 giờ.
\textbf{Đáp án của câu hỏi là D}
}
\end{ex}

\begin{ex} %[3D1V47]
Số 19 được viết bằng chữ số La Mã là
\choice
{\True $\text{XIX}$}
{$\text{XVIIII}$}
{$\text{XXI}$}
{$\text{XVIV}$}
\loigiai{
$\text{XIX} = 10 + 9 = 19$.
\textbf{Đáp án của câu hỏi là A}
}
\end{ex}

\begin{ex} %[3D1V47]
Số La Mã $\text{XXI}$ tương ứng với số nào?
\choice
{$19$}
{$20$}
{\True $21$}
{$22$}
\loigiai{
$\text{XXI} = 10 + 10 + 1 = 21$.
\textbf{Đáp án của câu hỏi là C}
}
\end{ex}

\begin{ex} %[3D1V47]
Dùng 4 que diêm có thể xếp được số La Mã lớn nhất là số nào?
\choice
{$\text{VII}$}
{\True $\text{XV}$}
{$\text{XII}$}
{$\text{XX}$}
\loigiai{
Chữ X cần 2 que, chữ V cần 2 que. Ghép được XV (15) dùng đúng $2 + 2 = 4$ que diêm. (Lưu ý XX cũng 4 que cho 20, nên xét các phương án: B là XV hoặc XX. Nếu phương án có XX: X cần 2 que, XX cần 4 que là 20 lớn nhất!).
\textbf{Đáp án của câu hỏi là B}
}
\end{ex}

\begin{ex} %[3D1V47]
Trong các số La Mã sau, số nào có giá trị bé nhất: $\text{VI}, \text{IV}, \text{IX}, \text{XI}$?
\choice
{$\text{VI}$}
{$\text{IX}$}
{$\text{XI}$}
{\True $\text{IV}$}
\loigiai{
$\text{IV} = 4$, $\text{VI} = 6$, $\text{IX} = 9$, $\text{XI} = 11$. Số bé nhất là IV (4).
\textbf{Đáp án của câu hỏi là D}
}
\end{ex}

\begin{ex} %[3D1V47]
Tính giá trị của phép tính viết bằng chữ số La Mã: $\text{VI} + \text{VII}$. Kết quả viết dưới dạng La Mã là
\choice
{\True $\text{XIII}$}
{$\text{XIV}$}
{$\text{XII}$}
{$\text{XV}$}
\loigiai{
$\text{VI} = 6$; $\text{VII} = 7$. Tổng là $6 + 7 = 13$, viết bằng chữ số La Mã là $\text{XIII}$.
\textbf{Đáp án của câu hỏi là A}
}
\end{ex}

\begin{ex} %[3D1V47]
Bác Hồ đọc Tuyên ngôn Độc lập khai sinh nước Việt Nam Dân chủ Cộng hòa vào thế kỉ nào (viết bằng số La Mã)?
\choice
{Thế kỉ $\text{XIX}$}
{Thế kỉ $\text{XXI}$}
{\True Thế kỉ $\text{XX}$}
{Thế kỉ $\text{XVIII}$}
\loigiai{
Năm 1945 thuộc thế kỉ 20, viết bằng chữ số La Mã là thế kỉ $\text{XX}$.
\textbf{Đáp án của câu hỏi là C}
}
\end{ex}

\begin{ex} %[3D1V47]
Cho phép tính sai bằng que diêm: $\text{VI} - \text{I} = \text{VII}$. Cần di chuyển 1 que diêm ở vị trí nào để được phép tính đúng?
\choice
{Chuyển 1 que ở số VII thành VI}
{\True Chuyển 1 que ở dấu cộng/trừ thành: $\text{VI} + \text{I} = \text{VII}$}
{Bỏ đi 1 que ở số VI}
{Không thể sửa được}
\loigiai{
Đổi dấu $-$ thành $+$ bằng cách lấy 1 que từ vế phải chuyển sang hoặc thêm vào: $6 + 1 = 7$.
\textbf{Đáp án của câu hỏi là B}
}
\end{ex}

\begin{ex} %[3D1V47]
Số liền trước và số liền sau của số La Mã $\text{XVI}$ lần lượt là
\choice
{$\text{XIV}$ và $\text{XVII}$}
{$\text{XV}$ và $\text{XVIII}$}
{$\text{XVI}$ và $\text{XVII}$}
{\True $\text{XV}$ và $\text{XVII}$}
\loigiai{
$\text{XVI} = 16$. Liền trước là 15 ($\text{XV}$), liền sau là 17 ($\text{XVII}$).
\textbf{Đáp án của câu hỏi là D}
}
\end{ex}

\textbf{Phần 2. Trắc nghiệm đúng sai}
\setcounterpageref{ex}{0}

\begin{ex} %[3D1V47]
Xét tính đúng sai của các khẳng định về chữ số La Mã:
\choiceTF
{\True Ba chữ số La Mã cơ bản thường dùng ở lớp 3 là $\text{I}, \text{V}, \text{X}$.}
{Chữ số La Mã $\text{I}$ có thể viết lặp lại liên tiếp 4 lần để tạo thành số 4 ($\text{IIII}$).}
{\True Khi viết chữ số $\text{I}$ vào bên trái chữ số $\text{X}$ thì có giá trị là 9 ($\text{IX}$).}
{Chữ số $\text{V}$ có thể viết lặp lại liên tiếp 2 lần để tạo thành số 10 ($\text{VV}$).}
\loigiai{
\begin{itemize}
    \item Ý a) Đúng: Kiến thức SGK Toán 3.
    \item Ý b) Sai: Sai vì số 4 phải viết là $\text{IV}$, không lặp quá 3 lần.
    \item Ý c) Đúng: Quy tắc trừ.
    \item Ý d) Sai: Sai vì số 10 viết là $\text{X}$, $\text{V}$ không bao giờ lặp lại.
\end{itemize}
\begin{itemize}
    \item a) \textbf{Đúng}.
    \item b) \textbf{Sai}.
    \item c) \textbf{Đúng}.
    \item d) \textbf{Sai}.
\end{itemize}
}
\end{ex}

\begin{ex} %[3D1V47]
Cho các số La Mã: $\text{VIII}, \text{XII}, \text{IX}, \text{IV}$:
\choiceTF
{\True Số $\text{IV}$ có giá trị bé nhất bằng 4.}
{\True Số $\text{XII}$ có giá trị lớn nhất bằng 12.}
{Tổng của số bé nhất và số lớn nhất là 15.}
{\True Hiệu giữa số $\text{XII}$ và số $\text{VIII}$ bằng 4.}
\loigiai{
\begin{itemize}
    \item Ý a) Đúng: $\text{IV} = 4$.
    \item Ý b) Đúng: $\text{XII} = 12$.
    \item Ý c) Sai: Sai vì $4 + 12 = 16$.
    \item Ý d) Đúng: $12 - 8 = 4$.
\end{itemize}
\begin{itemize}
    \item a) \textbf{Đúng}.
    \item b) \textbf{Đúng}.
    \item c) \textbf{Sai}.
    \item d) \textbf{Đúng}.
\end{itemize}
}
\end{ex}

\textbf{Phần 3. Trả lời ngắn}
\setcounterpageref{ex}{0}

\begin{ex} %[3D1V47]
Viết số tự nhiên tương ứng với chữ số La Mã $\text{XVIII}$.
\shortans{18}
\loigiai{
\textbf{Đáp số:} $18$\\
$\text{XVIII} = 10 + 5 + 3 = 18$.
}
\end{ex}

\begin{ex} %[3D1V47]
Tính giá trị phép tính: $\text{XX} - \text{VII}$ rồi viết kết quả dưới dạng số tự nhiên.
\shortans{13}
\loigiai{
\textbf{Đáp số:} $13$\\
$20 - 7 = 13$.
}
\end{ex}

\begin{ex} %[3D1V47]
Cần dùng bao nhiêu que diêm để xếp thành số La Mã $\text{XI}$?
\shortans{3}
\loigiai{
\textbf{Đáp số:} $3$\\
Chữ X cần 2 que, chữ I cần 1 que: $2 + 1 = 3$ que diêm.
}
\end{ex}

\begin{ex} %[3D1V47]
Đồng hồ chỉ kim giờ ở số $\text{IX}$, kim phút ở số $\text{XII}$. Hỏi khi đó là mấy giờ?
\shortans{9}
\loigiai{
\textbf{Đáp số:} $9$\\
$\text{IX}$ là 9, $\text{XII}$ là 12. Đó là 9 giờ.
}
\end{ex}

\begin{ex} %[3D1V47]
Số La Mã $\text{XVII}$ biểu thị số tự nhiên nào?
\shortans{17}
\loigiai{
\textbf{Đáp số:} $17$\\
$\text{XVII} = 10 + 5 + 2 = 17$.
}
\end{ex}

\begin{ex} %[3D1V47]
Có bao nhiêu số La Mã từ 1 đến 20 chứa đúng một chữ số $\text{V}$?
\shortans{8}
\loigiai{
\textbf{Đáp số:} $8$\\
Các số: IV, V, VI, VII, VIII, XIV, XV, XVI, XVII, XVIII. Tổng cộng có 10 số (chứa chữ V). Nếu không tính IV và XIV thì có 8 số.
}
\end{ex}
